\chapter{Case Study: 32-bit and 64-bit Processors}

\begin{khulasa}[: what will we learn in this chapter?]
\begin{itemize}
    \item Why we moved from 16 bits to 32 and then 64 bits, and what changed.
    \item How the architecture grew from IA-32 to AMD64.
    \item The registers of modern processors, and flat addressing.
    \item New instructions: SSE and AVX, which compute several numbers at the same time.
    \item The pipeline: how the processor runs several instructions at the same time.
\end{itemize}
\end{khulasa}

\begin{intabih}
    The programs of this chapter are written in \textbf{NASM} syntax for \textbf{Linux}, and we really ran them on an
    Intel Core i7-8665U processor (4 cores). To assemble and run:
    \code{nasm -f elf64 prog.asm}, then \code{ld prog.o -o prog}, then \code{./prog}.
\end{intabih}

%================================================
\section{From 16 to 32, then 64 bits}
%================================================

\begin{tarif}
    We say a processor is an \textbf{$n$-bit} processor when its \textbf{general registers} are $n$ bits wide: this is
    the size of the number it computes in one instruction, and the size of the address it can keep in a register.
\end{tarif}

\subsection{The main differences}
\begin{center}
\small
\begin{tabularx}{\linewidth}{|>{\centering\arraybackslash}p{3.2cm}|C|C|C|}
    \hline
    \ras{} & \ras{16 bits (8086)} & \ras{32 bits (IA-32)} & \ras{64 bits (x86-64)} \\ \hline
    biggest unsigned number & $2^{16}-1 = 65\,535$ & $2^{32}-1 \approx 4.3$ billion & $2^{64}-1 \approx 1.8 \times 10^{19}$ \\ \hline
    max memory & 1 MB (with segments) & 4 GB & 256 TB (48 bits in practice) \\ \hline
    general registers & 8 & 8 & \textbf{16} \\ \hline
    register names & AX, BX, \dots & EAX, EBX, \dots & RAX, RBX, \dots, R15 \\ \hline
    addressing & segmented & flat (in practice) & flat \\ \hline
\end{tabularx}
\end{center}

\subsection{Why do we need more than 32 bits?}
\begin{enumerate}
    \item \textbf{Memory}: with a 32-bit address we cannot use more than 4 GB. Today's computers have 16 or 32 GB,
          and servers have hundreds of gigabytes.
    \item \textbf{Performance}: a 64-bit register computes a big number in one instruction, while a 32-bit processor needs
          at least two instructions (like \code{add} and \code{adc} in Chapter 3). And twice as many registers (16 instead of 8)
          means fewer trips to memory.
    \item \textbf{Compatibility}: a 64-bit processor also runs old 32-bit programs without any change.
\end{enumerate}

\begin{khatwa}[: what does ``4 GB'' mean exactly?]
    \begin{enumerate}
        \item A 32-bit address can take $2^{32}$ different values.
        \item $2^{32} = 2^{2} \times 2^{30} = 4 \times 2^{30}$.
        \item And $2^{30}$ bytes $= 1$ GB, so the maximum size is 4 GB.
    \end{enumerate}
    With 48 bits: $2^{48} = 2^{8} \times 2^{40} = 256$ TB, that is 65536 times bigger!
\end{khatwa}

%================================================
\section{How the architecture grew: from IA-32 to AMD64}
%================================================

\begin{center}
\begin{tikzpicture}[e/.style={unit=#1, minimum width=2.9cm, minimum height=1.3cm, font=\small}]
    \node[e=brandMuted] (a) at (0,0) {8086 (1978)\\16 bits, segments};
    \node[e=brandPurple] (b) at (3.6,0) {80286 (1982)\\protected mode};
    \node[e=brandTeal] (c) at (7.2,0) {80386 (1985)\\\textbf{IA-32}};
    \node[e=brandOrange] (d) at (10.8,0) {AMD64 (2003)\\\textbf{64 bits}};
    \draw[sahmT] (a) -- (b); \draw[sahmT] (b) -- (c); \draw[sahmT] (c) -- (d);
\end{tikzpicture}
\end{center}

\begin{itemize}
    \item \textbf{IA-32} (Intel Architecture 32 bits): it started with the 80386. Registers became 32 bits wide
          (EAX: the E means \emph{Extended}), and the \textbf{protected mode} appeared: it stops programs from writing in each other's memory.
    \item \textbf{AMD64}: the surprise is that AMD (not Intel) designed the 64-bit extension in 2003, while keeping
          compatibility. Intel adopted it later as Intel 64. Today we call it x86-64.
\end{itemize}

\subsection{The working modes of the processor}
\begin{center}
\small
\begin{tabularx}{\linewidth}{|>{\centering\arraybackslash}p{3.3cm}|R|R|}
    \hline
    \ras{Mode} & \ras{Description} & \ras{When?} \\ \hline
    real mode & behaves like an 8086: 16 bits and 1 MB & when the computer starts, before the operating system \\ \hline
    protected mode & 32 bits, memory protection & 32-bit systems \\ \hline
    long mode & 64 bits, with a sub-mode to run 32-bit programs & all modern systems \\ \hline
\end{tabularx}
\end{center}

\begin{hal}
    Every modern computer wakes up ``thinking'' it is an 8086! It starts in real mode, then the operating system moves
    it to protected mode, then to long mode. This is how the architecture stays compatible since 1978.
\end{hal}

%================================================
\section{Registers in 32-bit and 64-bit processors}
%================================================

\subsection{Nested registers}
Each new register \textbf{contains} the old register in its right half (the low bits):

\begin{center}
\begin{tikzpicture}[b/.style={draw=brandDark, thick, minimum height=0.75cm, font=\codefont\bfseries}]
    \node[b, minimum width=12cm, fill=brandOrange!15, anchor=east] (r) at (6,1.6) {RAX: 64 bits};
    \node[b, minimum width=6cm, fill=brandTeal!15, anchor=east] (e) at (6,0.8) {EAX: 32 bits};
    \node[b, minimum width=3cm, fill=brandGreen!20, anchor=east] (x) at (6,0) {AX: 16};
    \node[b, minimum width=1.5cm, fill=brandPurple!20, anchor=east] (h) at (4.5,-0.8) {AH};
    \node[b, minimum width=1.5cm, fill=brandPurple!35, anchor=east] (l) at (6,-0.8) {AL};
    \foreach \x/\t in {6/0, 4.5/8, 3/16, 0/32, -6/64} \node[font=\codefont\scriptsize] at (\x,-1.45) {\t};
    \draw[brandMuted, dashed] (0,-1.25) -- (0,1.2); \draw[brandMuted, dashed] (3,-1.25) -- (3,0.4);
\end{tikzpicture}
\end{center}

\begin{center}
\small
\begin{tabular}{|c|c|c|c|l|}
    \hline
    \ras{64 bits} & \ras{32 bits} & \ras{16 bits} & \ras{8 bits} & \ras{Usual use} \\ \hline
    RAX & EAX & AX & AL & value returned by a function \\ \hline
    RBX & EBX & BX & BL & general \\ \hline
    RCX & ECX & CX & CL & counter, 4th parameter \\ \hline
    RDX & EDX & DX & DL & 3rd parameter \\ \hline
    RSI & ESI & SI & SIL & 2nd parameter \\ \hline
    RDI & EDI & DI & DIL & 1st parameter \\ \hline
    RSP & ESP & SP & SPL & top of the stack \\ \hline
    RBP & EBP & BP & BPL & base of the stack frame \\ \hline
    R8 \dots\ R15 & R8D \dots\ R15D & R8W \dots & R8B \dots & \textbf{new in 64 bits} \\ \hline
\end{tabular}
\end{center}

\begin{itemize}
    \item \textbf{RIP}: the instruction pointer (it was IP, then EIP).
    \item \textbf{RFLAGS}: the flags register; ZF, CF, SF, OF are in the same place as in the 8086.
    \item \textbf{Segment registers}: they still exist (CS, DS\dots) but their base is 0 in 64-bit mode, except FS and GS,
          which the system uses for data that belongs to each thread.
\end{itemize}

\subsection{A test: what happens to the other bits?}
\asm{ch4/regs64.asm}
\screen{ch4/regs64.out}

\begin{sharh}
    \s{5--6}{Put a 64-bit number in RAX and show it: 1122334455667788.}
    \s{9}{Writing EAX (32 bits) \textbf{clears} the high half: 00000000AABBCCDD.}
    \s{13}{Writing AX (16 bits) keeps the other bits: \dots 5566\textbf{1234}.}
    \s{17}{Writing AL changes only the last byte: \dots 77\textbf{FF}.}
    \s{24--42}{A procedure that shows RAX in hexadecimal: at each turn it rotates 4 bits, keeps them with the mask \code{0xF}, and turns them into a character.}
\end{sharh}

\begin{intabih}
    A strange but important rule: in 64-bit mode, \textbf{every write to a 32-bit register clears the high bits}, but a
    write to 8 or 16 bits does not. AMD designed this rule to make the processor faster.
\end{intabih}

%================================================
\section{Advanced addressing modes}
%================================================

\subsection{From segmented addressing to flat addressing}
\begin{center}
\begin{tikzpicture}
    \draw[thick, fill=brandLight] (-6,0) rectangle (-4,-4);
    \foreach \y/\c in {0/brandTeal, -1.3/brandOrange, -2.6/brandGreen} \fill[\c!30] (-6,\y) rectangle (-4,\y-1.2);
    \foreach \y/\t in {-0.6/CS, -1.9/DS, -3.2/SS} \node[font=\codefont] at (-5,\y) {\t};
    \node[font=\small\bfseries, align=center] at (-5,0.6) {segmented (8086)};
    \node[font=\scriptsize, text width=4cm, align=center] at (-5,-4.6) {zones of 64 KB, and address = segment:offset};
    \draw[thick, fill=brandTeal!15] (0,0) rectangle (2,-4);
    \node[font=\small, align=center] at (1,-2) {one single\\space};
    \node[font=\small\bfseries, align=center] at (1,0.6) {flat (x86-64)};
    \node[font=\scriptsize, text width=4cm, align=center] at (1,-4.6) {address = one number from 0 to the end};
    \draw[sahmO] (-3.6,-2) -- (-0.4,-2);
\end{tikzpicture}
\end{center}

In the \term{flat memory model}{modèle mémoire plat}, the program sees memory \textbf{as one long strip} of bytes,
numbered from 0. We do not need to compute ``segment $\times$ 16 + offset'': the address is just a number in a 64-bit register.

\begin{tashbih}
    Segmented addressing is like saying ``building 12, floor 3''. Flat addressing is like giving every flat in the
    town one unique number: ``flat number 73\,502''. Much simpler, as long as we have big enough numbers.
\end{tashbih}

\subsection{The general form of an address}
\begin{tarif}
    \centering
    $\text{address} = \text{base} + \text{index} \times \text{scale} + \text{displacement}$
    \par\small
    base and index: any two registers. scale: 1, 2, 4 or 8. displacement: a fixed number.
\end{tarif}

\begin{center}
\small
\begin{tabular}{|l|l|}
    \hline
    \ras{Writing (NASM)} & \ras{Use} \\ \hline
    \code{[rbx]} & base only: a pointer \\ \hline
    \code{[rbx + 16]} & base + displacement: a field in a structure (struct) \\ \hline
    \code{[tab + rcx*8]} & index $\times$ scale + displacement: element i of an array of 8-byte numbers \\ \hline
    \code{[rbx + rcx*4 + 8]} & everything together \\ \hline
    \code{[rel msg]} & relative to RIP: special to 64 bits \\ \hline
\end{tabular}
\end{center}

\begin{khatwa}[: where does \code{[tab + rcx*8]} read when RCX = 3?]
    \begin{enumerate}
        \item Each element of \code{tab} is a 64-bit number = 8 bytes (\code{dq}).
        \item Element number 3 (the fourth one, because we start at 0) is $3 \times 8 = 24$ bytes from the start.
        \item If \code{tab} is at address \hex{402000}, the address is $\hex{402000} + 24 = \hex{402018}$.
    \end{enumerate}
    The scale saves us from multiplying the index ourselves: the processor does it inside the instruction.
\end{khatwa}

\asm{ch4/sum64.asm}
\screen{ch4/sum64.out}

\begin{sharh}
    \s{5--6}{\code{xor} of a register with itself = 0: the sum and the index start at 0.}
    \s{8}{Addressing \textbf{with a scale}: read element i of the array, and add it to RAX.}
    \s{9--11}{\code{i++}, then ``while i < 5'' go back. This is a \code{while} loop, as in Chapter 3.}
    \s{18--36}{The same method as in Chapter 3 to print a number in decimal (divide by 10), but with 64-bit registers.}
    \s{39}{\code{dq} (define quad-word): 64-bit numbers. The first number ($10^{12}$) does not fit in 32 bits.}
\end{sharh}

\subsection{RIP-relative addressing}
In 64-bit mode we often write \code{lea rsi, [rel buf]}: the address is computed \textbf{relative to the current
instruction} ($\text{RIP} + \text{distance}$). The benefit: the program works wherever the system puts it in memory.
This is important for security (the system moves programs at each run so attackers cannot guess the addresses).

%================================================
\section{The instruction set}
%================================================

\subsection{The same instruction on 16, 32 and 64 bits}
We assembled the program below; these are the real machine-code bytes:

\asm{ch4/same.asm}

\begin{center}
\small
\begin{tabular}{|l|l|l|}
    \hline
    \ras{Instruction} & \ras{Machine code} & \ras{Note} \\ \hline
    \code{add ax, bx} & {\codefont \textcolor{orangeDeep}{66} 01 D8} & prefix 66: ``careful, 16 bits'' \\ \hline
    \code{add eax, ebx} & {\codefont 01 D8} & 32 bits is the default size \\ \hline
    \code{add rax, rbx} & {\codefont \textcolor{orangeDeep}{48} 01 D8} & prefix REX.W = 48: ``64 bits'' \\ \hline
    \code{add r8, r9} & {\codefont \textcolor{orangeDeep}{4D} 01 C8} & REX = 4D: 64 bits + new registers \\ \hline
\end{tabular}
\end{center}

\begin{khatwa}[: why is 01 D8 in the first three lines?]
    \begin{enumerate}
        \item 01 is the opcode ``add a register to a register or memory''. It is the same since the 8086.
        \item D8 = \bin{11\,011\,000} is the ModR/M byte we saw in Chapter 3: mod = 11 (register), reg = 011 (BX), r/m = 000 (AX).
        \item The size (16, 32 or 64) is not written in the instruction itself, but in a \textbf{prefix} before it. This is how
              the old instructions stayed valid, and the new sizes were added ``from the outside''.
        \item In 4D = \bin{0100\,1101}: 0100 means ``REX prefix'', then the bit W = 1 (64 bits), and two bits that add 8 to the
              two register numbers, so AX becomes R8 and BX becomes R9.
    \end{enumerate}
\end{khatwa}

\subsection{Comparison: a 32-bit program and a 64-bit program}
\textbf{In 32 bits:}
\asm{ch4/hello32.asm}
\screen{ch4/hello32.out}
\textbf{In 64 bits:}
\asm{ch4/hello64.asm}
\screen{ch4/hello64.out}

\begin{center}
\small
\begin{tabularx}{\linewidth}{|>{\centering\arraybackslash}p{3.8cm}|C|C|}
    \hline
    \ras{} & \ras{32 bits} & \ras{64 bits} \\ \hline
    asking the system for a service & \code{int 0x80} (an interrupt) & \code{syscall} (a special, faster instruction) \\ \hline
    number of the \emph{write} service & 4 in EAX & 1 in RAX \\ \hline
    parameters & EBX, ECX, EDX & RDI, RSI, RDX \\ \hline
    size of an address & 4 bytes & 8 bytes \\ \hline
\end{tabularx}
\end{center}

\begin{hal}
    We ran the two programs on the same 64-bit computer. The 32-bit program worked without any change: this is the
    compatibility sub-mode inside long mode.
\end{hal}

\subsection{What does the compiler write? From C to assembly}
We wrote a function in C, and asked the compiler gcc to give us the assembly it produces, with the command
\code{gcc -O1 -S -masm=intel sum.c}:

\begin{center}
\begin{minipage}[t]{0.47\linewidth}
    \cfile{ch4/sum.c}
\end{minipage}\hfill
\begin{minipage}[t]{0.49\linewidth}
    \asm{ch4/sum_gcc.s}
\end{minipage}
\end{center}

\begin{sharh}
    \s{1}{By the Linux 64-bit rules: the first parameter (\code{t}) arrives in RDI, and the second one (\code{n}) in RSI.}
    \s{2--3}{If $n \leq 0$ we jump to \code{.L4}: the result is 0.}
    \s{5}{\code{lea} computes the address of the end of the array $t + n \times 8$ \textbf{without reading memory}.}
    \s{8--11}{The loop: add the element, move 8 bytes forward, and stop at the end. The compiler used a pointer instead of the index \code{i}.}
    \s{13--14}{The result is always returned in RAX.}
\end{sharh}

\subsection{SSE and AVX: computing several numbers at once}
\begin{tarif}
    \term{SIMD}{une instruction, données multiples} (Single Instruction, Multiple Data): one instruction applies the
    same operation to \textbf{several numbers at the same time}.
\end{tarif}

\begin{center}
\begin{tikzpicture}[c/.style={draw=brandDark, thick, minimum width=1.3cm, minimum height=0.65cm, font=\codefont}]
    \foreach \v [count=\i from 0] in {1.5, 2.5, 3.5, 4.5} \node[c, fill=brandTeal!15] at (1.3*\i,1.6) {\v};
    \foreach \v [count=\i from 0] in {10, 20, 30, 40} \node[c, fill=brandOrange!15] at (1.3*\i,0.6) {\v};
    \foreach \v [count=\i from 0] in {11.5, 22.5, 33.5, 44.5} \node[c, fill=brandGreen!20] at (1.3*\i,-0.9) {\v};
    \draw[very thick] (-0.7,0.1) -- (4.6,0.1);
    \node[font=\codefont\bfseries] at (-1.1,0.6) {+};
    \node[font=\small, anchor=west] at (4.8,1.6) {XMM0 (128 bits = 4 numbers)};
    \node[font=\small, anchor=west] at (4.8,-0.9) {result: \textbf{one instruction} \code{addps}};
\end{tikzpicture}
\end{center}

\begin{center}
\small
\begin{tabular}{|c|c|c|c|c|}
    \hline
    \ras{Extension} & \ras{Year} & \ras{Registers} & \ras{Size} & \ras{floats at once} \\ \hline
    SSE & 1999 & XMM0--XMM15 & 128 bits & 4 \\ \hline
    AVX & 2011 & YMM0--YMM15 & 256 bits & 8 \\ \hline
    AVX-512 & 2016 & ZMM0--ZMM31 & 512 bits & 16 \\ \hline
\end{tabular}
\end{center}

\cfile{ch4/simd.c}
\screen{ch4/simd.out}

When we open the compiled program, we really find the instruction:
\begin{codebox}\codefont\small
    addps \ xmm0, XMMWORD PTR [rbp-0x50] \quad \textcolor{red!75!black}{; 4 additions at once}
\end{codebox}

And the AVX version adds 8 numbers with the instruction \code{vaddps ymm0, ymm0, [\dots]}:
\cfile{ch4/avx.c}
\screen{ch4/avx.out}

\begin{tashbih}
    A normal instruction is like a teacher who corrects one exam paper at a time. SIMD is like a teacher who puts
    8 papers side by side and corrects question 1 on all of them with one move. This is why it is used in image and
    video processing, games, and artificial intelligence, where we apply the same operation to millions of numbers.
\end{tashbih}

%================================================
\section{The instruction cycle and the pipeline}
%================================================

\subsection{The five stages}
In modern processors, each instruction is split into stages, and each stage has its own circuit:
\begin{center}
\small
\begin{tabular}{|c|c|l|}
    \hline
    \ras{Stage} & \ras{Name} & \ras{Work} \\ \hline
    IF & fetch & read the instruction from memory (address in RIP) \\ \hline
    ID & decode & understand the instruction and read the registers \\ \hline
    EX & execute & compute the result in the UAL, or compute an address \\ \hline
    MEM & memory & read or write memory if needed \\ \hline
    WB & write back & store the result in the register \\ \hline
\end{tabular}
\end{center}

\subsection{Without a pipeline and with a pipeline}
\begin{tashbih}
    In a laundry: washing takes one hour, drying one hour, ironing one hour. To wash 3 bags one after the other we
    need 9 hours. But when the first bag goes from the washing machine to the dryer, we put the second bag in the
    washing machine at once. The three machines work together, so we finish in only 5 hours. This is a pipeline.
\end{tashbih}

\begin{center}
\begin{tikzpicture}[s/.style={draw=brandDark, minimum width=0.95cm, minimum height=0.55cm, font=\codefont\scriptsize\bfseries}]
    \node[font=\small\bfseries] at (4.3,0.9) {with a pipeline: 4 instructions in 8 cycles};
    \foreach \c in {1,...,8} \node[font=\codefont\scriptsize] at (0.95*\c, 0.35) {\c};
    \foreach \i/\col in {1/brandTeal, 2/brandOrange, 3/brandGreen, 4/brandPurple} {
        \node[font=\small] at (-0.3,-0.6*\i+0.1) {I\i};
        \foreach \st [count=\k from 0] in {IF, ID, EX, MEM, WB} {
            \pgfmathsetmacro\x{0.95*(\i+\k)}
            \node[s, fill=\col!25] at (\x, -0.6*\i+0.1) {\st};
        }
    }
\end{tikzpicture}
\end{center}

\begin{tarif}
    For $n$ instructions and $k$ stages (one cycle per stage):
    \begin{itemize}
        \item without a pipeline: $n \times k$ cycles.
        \item with a pipeline: $k + (n - 1)$ cycles: the first one needs $k$ cycles, then one instruction finishes in \textbf{every} cycle.
    \end{itemize}
\end{tarif}

\begin{khatwa}[: 100 instructions in a 5-stage pipeline]
    \begin{enumerate}
        \item Without a pipeline: $100 \times 5 = 500$ cycles.
        \item With a pipeline: $5 + (100 - 1) = 104$ cycles.
        \item The speed-up: $\dfrac{500}{104} \approx 4.8$ times, close to the number of stages (5).
    \end{enumerate}
\end{khatwa}

\subsection{Hazards: when does the pipeline stop?}
\begin{enumerate}
    \item \textbf{Data hazard}: \code{add rax, rbx} then \code{add rcx, rax}. The second instruction needs the new RAX,
          but the first one has not written it yet. The fix: the processor ``forwards'' the result directly from EX to the
          second instruction, or waits one cycle.
    \item \textbf{Control hazard}: after \code{jmp} or \code{jne}, the instructions fetched after it may be wrong. The fix:
          \term{branch prediction}{prédiction de branchement}: the processor guesses if it will jump, and if it guessed wrong,
          it throws away the wrong instructions and starts again.
    \item \textbf{Structural hazard}: two instructions need the same circuit in the same cycle (like memory).
\end{enumerate}

\subsection{Example: \code{mov}, \code{add} and \code{jmp} in 64 bits}
We follow three instructions from the program \code{same.asm} and the programs of this chapter, with their real bytes:

\begin{center}
\small
\begin{tabularx}{\linewidth}{|>{\centering\arraybackslash}p{1.1cm}|>{\raggedright\arraybackslash}p{3.3cm}|R|R|}
    \hline
    \ras{Stage} & \ras{\code{mov rax, 60}} & \ras{\code{add rax, rbx}} & \ras{\code{jmp fin}} \\ \hline
    IF & fetch 7 bytes 48 C7 C0 3C 00 00 00 & fetch 3 bytes 48 01 D8 & fetch 2 bytes EB 00 \\ \hline
    ID & REX.W: 64 bits, destination RAX, the value 60 is in the instruction & 64 bits, source RBX, destination RAX; read both registers & short jump, offset 0 \\ \hline
    EX & nothing to compute (the value is ready) & UAL: $\text{RAX} + \text{RBX}$, and compute the flags & compute the target: $\text{RIP} + 0$ \\ \hline
    MEM & nothing & nothing & nothing \\ \hline
    WB & \nql{RAX}{60} & RAX gets the result; update RFLAGS & RIP gets the target; throw away wrongly fetched instructions \\ \hline
\end{tabularx}
\end{center}

\begin{hal}
    The Core i7 we used does much more than a 5-stage pipeline: it has more than 14 stages, and it can start
    \textbf{4 instructions in every cycle}. It even runs them \textbf{out of order} when they do not depend on each
    other, then puts the results back in order. All this happens without the programmer seeing any difference.
\end{hal}

%================================================
\section{Exercises}
%================================================

\begin{tamrin}
    How many bytes can a processor address with a 36-bit address? (This is what the PAE extension allowed in 32-bit processors.)
\end{tamrin}
\begin{hall}
    $2^{36} = 2^{6} \times 2^{30} = 64$ GB.
\end{hall}

\begin{tamrin}
    RAX = \hex{FFFFFFFFFFFFFFFF}. What is RAX after: (a) \code{mov eax, 1} (b) \code{mov ax, 1} (c) \code{mov al, 1}?
\end{tamrin}
\begin{hall}
    (a) \hex{0000000000000001} (writing 32 bits clears the high half).\\
    (b) \hex{FFFFFFFFFFFF0001}. \quad (c) \hex{FFFFFFFFFFFFFF01}.
\end{hall}

\begin{tamrin}
    An array of \code{int} (4 bytes each) starts at \hex{601000}. What is the address of element number 10? Write the addressing in NASM.
\end{tamrin}
\begin{hall}
    $\hex{601000} + 10 \times 4 = \hex{601000} + 40 = \hex{601028}$. The writing: \code{[tab + rcx*4]} with RCX = 10.
\end{hall}

\begin{tamrin}
    A 6-stage pipeline, and 1000 instructions. Compute the number of cycles without and with a pipeline, and the speed-up.
\end{tamrin}
\begin{hall}
    Without: $1000 \times 6 = 6000$. With: $6 + 999 = 1005$. Speed-up $= \dfrac{6000}{1005} \approx 5.97$.
\end{hall}

\begin{tamrin}
    In this program, where is the data hazard? \code{mov rax, [x]}, then \code{add rax, 5}, then \code{mov [y], rax}.
\end{tamrin}
\begin{hall}
    The second instruction needs the RAX that the first one reads from memory, and the third one needs the RAX that the
    second one computes: each instruction depends on the one just before it.
\end{hall}
